%-----------PROJECTS-----------
\section{Projects}
  \resumeSubHeadingListStart

% --- Kalapa Behavior Score Flutter App ---
\resumeProjectHeading
  {Kalapa Behavior Score – Flutter Credit Scoring Demo App}
  {Flutter, Riverpod (Generator), go\_router, Freezed, Dio, MethodChannel}
\resumeItemListStart
  \resumeItem{Built a Flutter client showcasing the \textbf{Kalapa Behavior Score SDK} in a simulated consumer-loan flow, using \textbf{Feature-First Clean Architecture} (presentation/domain/data/core) for real-time on-device credit/risk scoring.}
  \resumeItem{Used \textbf{Riverpod} (Generator) for state management/DI, \textbf{go\_router} for navigation, and \textbf{Freezed}/\textbf{fpdart} for immutable models and functional error handling.}
  \resumeItem{Bridged to the native Android Behavior Score SDK via \textbf{MethodChannel}, with a mock data source fallback for iOS/desktop, and implemented permission flows for on-device telemetry collection.}
  \resumeItem{Shipped dev/prod build flavors with obfuscated, minified release builds (R8, split debug symbols) for Google Play distribution.}
\resumeItemListEnd

% --- LOS Miniapp Host ---
\resumeProjectHeading
  {LOS Miniapp Host – Flutter WebView Loan-Registration App}
  {Flutter, Riverpod (Generator), go\_router, InAppWebView, Dio, Freezed}
\resumeItemListStart
  \resumeItem{Built a Flutter host app embedding a partner mini-app (FE CREDIT loan registration) via an in-app \textbf{WebView} instead of handing off to an external browser, on \textbf{Feature-First Clean Architecture} (presentation/domain/data/core).}
  \resumeItem{Used \textbf{Riverpod} (Generator, \texttt{@riverpod}) for compile-time-safe DI/state with \texttt{autoDispose} lifecycle management, and \textbf{go\_router} for declarative routing and deep linking.}
  \resumeItem{Implemented enterprise structured logging with redaction, \textbf{fpdart}-based functional error handling, and a Material 3 design system with WCAG-compliant theming and English/Vietnamese localization.}
\resumeItemListEnd

% --- Viettel Verification App ---
\resumeProjectHeading
  {Viettel EID/Passport Verification App – Android Kiosk Solution}
  {Kotlin, MVVM, Jetpack, ML Kit, REST API}
\resumeItemListStart
  \resumeItem{Developed an Android kiosk application in \textbf{Kotlin} using \textbf{MVVM + Clean Architecture}, deployed on Joyusing Z10S Pro devices with robust lifecycle handling and hardware integration.}
  \resumeItem{Integrated manufacturer SDKs for camera, NFC reader, and biometric modules, implementing asynchronous capture/processing flows with \textbf{Coroutines} for smooth and stable real-time operations.}
  \resumeItem{Implemented a full identity verification pipeline: document capture, MRZ OCR, ML Kit image preprocessing, NFC chip reading, and face-matching via Face Engine API.}
  \resumeItem{Connected with Viettel backend systems using \textbf{Retrofit} and secure REST APIs, and integrated \textbf{VNPAY sandbox QR payment} along with digital signing modules for transaction validation.}
\resumeItemListEnd

% --- eKYC Android SDK Modularization ---
\resumeProjectHeading
  {eKYC Android SDK – Multi-Module Architecture Refactor}
  {Kotlin, Gradle Multi-Module, Clean Architecture, CI/CD}
\resumeItemListStart
  \resumeItem{Refactored a monolithic Android eKYC SDK into a multi-module architecture (foundation, platform, capabilities, facades, composition), separating publishable host-facing facades (\textbf{document-sdk}, \textbf{liveness-sdk}, \textbf{nfc-sdk}) from the full \textbf{ekyc-sdk} composition root.}
  \resumeItem{Enforced a strict one-directional dependency graph (capability to facade to composition, no peer-capability edges) via a Gradle-level dependency matrix check to prevent architectural drift.}
  \resumeItem{Designed capability modules (camera, document capture, liveness detection, NFC) behind shared \textbf{core} contracts/gateways, plus a centralized localization system and a numbered error-code taxonomy (\textbf{KalapaSDKException}) for consistent host integration.}
\resumeItemListEnd

  \resumeSubHeadingListEnd
